Choose a fixed integral $F_2$ divisible over $A$ by the
characteristic elements of $\mathbf H^2_{\rm loc}(T\chi)$
for the finitely many allowed local types at two and finite
sign components. Its central value can be chosen nonzero:
derived local control, and the absence of local Iwasawa $H^3$,
identify its rational central $H^2$ with $H^2(\mathbf Q_2,V\chi)=0$.
The last equality follows from local duality and finiteness of
local elliptic torsion. Hence the characteristic elements have
no factor $t$. Kato (12.2.3), p.~220, gives the required finiteness
and torsion of this local Iwasawa module.

At the finitely many old bad primes, choose stable lattices in
the $W_q$ and multiply all the determinants $D_q$ for the finitely
many local twisting types. Clear their fixed scalar denominators.
Their central values are nonzero, since the cokernel of
$1-\operatorname{Fr}_{W_q}$ is $H^2(\mathbf Q_q,V\chi)=0$
by the same local duality. Call this product $F_{\rm bad}$ and set
$F_E=F_2F_{\rm bad}$.

The quotient $\mathbf H^2/\mathbf H^2_0$ injects into the local
$\mathbf H^2_{\rm loc}$, so $F_2$ pays for it in (K15).
Multiplying a class in a free rank-one module by a nonzero scalar
adds that scalar's valuation to its index. Formula (E.1) already
pays for every unused good determinant in (K16); $F_{\rm bad}$
pays for the remaining old ones. Adding these length inequalities
proves (K2).
Here is the length calculation with all contributions separated.
Fix a height-one prime $\mathfrak p$ of $A$ and write
$R=A_{\mathfrak p}$. Let $M$ be the localized degree-one module
for the primitive $j_*$ system, let $y=y_{2^\infty}$, and write
$\operatorname{ind}_R(y)=\ell_R(M/Ry)$.
The preceding freeness argument gives $M\simeq R$.
Put $Q_\chi=\prod_{q\mid F/d}E_q(\chi)$.
All primes in $F/d$ are good: the family has identical local
characters at primes dividing $N$, so a bad prime cannot ramify
in one constituent and be unramified in another.
The cone (K16) has zero degree-one cohomology, hence identifies
$M$ with the localized full degree-one module. Its long exact
sequence and (K15) give
\[
\begin{split}
 \ell_R H^2(C_{S,x})_{\mathfrak p}
 &\le \ell_R\mathbf H^2(T\chi)_{0,\mathfrak p}
       +\ell_R\mathbf H^2_{\rm loc}(T\chi)_{\mathfrak p}
       +\sum_{q\in S,\ q\ne2}v_{\mathfrak p}(D_q)\\
 &\le \operatorname{ind}_R(y)
       +v_{\mathfrak p}(F_2)+v_{\mathfrak p}(F_{\rm bad})
       +v_{\mathfrak p}(Q_\chi)\\
 &=\operatorname{ind}_R(F_EQ_\chi y)
  =\operatorname{ind}_R\bigl(F_E(Z_F)_x\bigr).
\end{split}
\]
The second line uses $D_q=1$ at active good primes, and
$D_q=E_q(\chi)$ at inactive good primes. It does not charge an
extra fixed constant for each active prime.

For completeness, nonvanishing at the center of the fixed
correction factors follows before multiplying over local types.
Local derived control and $H^3_{\rm Iw}=0$ identify
$\mathbf H^2_{\rm loc}/t\mathbf H^2_{\rm loc}$ rationally with
$H^2(\mathbf Q_2,V\chi)$. Local duality identifies the latter
with the dual of $H^0(\mathbf Q_2,V\chi)$, which is zero because
local elliptic torsion is finite. Nakayama at $(t)$ therefore
excludes a factor $t$ from its characteristic element.
At an old odd prime the cokernel of
$1-\operatorname{Fr}_{W_q}$ is $H^2(\mathbf Q_q,V\chi)=0$,
by the same duality. Its determinant at $t=0$ is nonzero.
Taking products over the finite list of local types, followed by
one fixed power of two to clear scalar denominators, gives
$F_E\in\Lambda$ with $F_E(0)\ne0$. This operation is independent
of the binary dimension and all later good conductor primes.


